\section{The global limiting Euler-kernel extremum}
\label{sec:kernel-extrema}

The pointwise kernel optimum and the optimum after averaging over
polylogarithm orders are different quantities. This section resolves the
limiting-kernel conjecture in the incoming sharp-Euler report
\cite{SharpEulerIncoming}. The argument also proves eventual uniqueness
of the maximizer and gives a convergent expansion in the reciprocal
truncation index.

For an integer $N\ge1$, put
\begin{equation}\label{eq:kernel-def}
 W_N(x)=\frac{(1-x)^N}{1+x},\qquad
 K_N(r,y)=\frac{W_N(ry^2)-W_N(r)}{1-y},\qquad
 M_N=\max_{(r,y)\in[0,1]^2}K_N(r,y).
\end{equation}
At $y=1$ the definition means $K_N(r,1)=-2rW_N'(r)$.
The divided difference is continuous on the closed square; hence the
maximum exists. We will use $c=Nr$ for the rescaled first coordinate.

\subsection{A decreasing envelope with uniform endpoint control}

\begin{lemma}[Exact power envelope]\label{lem:power-envelope}
For every integer $m\ge2$, define
\begin{equation}\label{eq:Bm}
 B_m(y)=\frac{(1-y^2)^m}
 {(1-y)(1-y^{2m/(m-1)})^{m-1}},\qquad 0<y<1.
\end{equation}
Its continuous endpoint values are
\[
 B_m(0)=1,\qquad B_m(1)=2\left(1-\frac1m\right)^{m-1}.
\]
Then
\begin{equation}\label{eq:power-max}
 B_m(y)=\max_{0\le r\le1}
 \frac{(1-ry^2)^m-(1-r)^m}{1-y}.
\end{equation}
The sequence $B_m$ decreases pointwise on $[0,1]$ to the continuous
function
\begin{equation}\label{eq:B-limit}
 B(y)=(1+y)y^{2y^2/(1-y^2)},\qquad B(0)=1,\quad B(1)=2/e.
\end{equation}
The convergence is uniform on $[0,1]$.
\end{lemma}
\begin{proof}
For $v=y^2\in(0,1)$, differentiation in $r$ shows that the unique
maximum of $(1-rv)^m-(1-r)^m$ satisfies
\[
 \frac{1-r}{1-rv}=v^{1/(m-1)}.
\]
Substitution proves \eqref{eq:power-max} and \eqref{eq:Bm}.
The endpoint $y=1$ is also the direct maximum of
$2mr(1-r)^{m-1}$.

To prove monotonicity, set $q=1/(m-1)$ and
$L(q)=\log(1-v^{1+q})$. Then
\[
 \log B_m(y)=\log(1+y)-\frac{L(q)-L(0)}q.
\]
The function $L$ is strictly concave, since
\[
 L''(q)=-\frac{(\log v)^2v^{1+q}}{(1-v^{1+q})^2}<0.
\]
Its secant slope from zero decreases as $q$ increases. Thus decreasing
$q$, which means increasing $m$, decreases $B_m(y)$.
As $q\downarrow0$, the secant slope tends to
$-v\log v/(1-v)$, giving \eqref{eq:B-limit}.
At $y=1$ the same monotonicity follows from the monotone increase of
$(1+1/k)^k$ to $e$, with $k=m-1$; at zero all values are one.
The limiting function is continuous at both endpoints because
$y^2\log y\to0$ at zero and
$2y^2\log y/(1-y^2)\to-1$ at one.
Dini's theorem now gives uniform convergence on the compact interval.
\end{proof}

\begin{lemma}[Uniform domination and local scaling]\label{lem:kernel-scale}
For $N\ge2$,
\begin{equation}\label{eq:kernel-envelope}
 K_N(r,y)\le B_N(y)\qquad(0\le r,y\le1).
\end{equation}
Moreover, on every compact rectangle $0\le c\le A$, $0\le y\le1$,
\begin{equation}\label{eq:kernel-limit}
 K_N(c/N,y)\longrightarrow
 H(c,y):=\frac{e^{-cy^2}-e^{-c}}{1-y},\qquad H(c,1)=2ce^{-c},
\end{equation}
uniformly, together with every fixed number of derivatives on compact
sets where the variables are interior.
\end{lemma}
\begin{proof}
The positive geometric expansion
\[
 W_N(x)=\sum_{k\ge0}2^{-k-1}(1-x)^{N+k}
\]
converges uniformly on $[0,1]$. Apply \cref{lem:power-envelope}
termwise and use $B_{N+k}\le B_N$. This also proves the continuous
endpoint case by passage to the limit.

For bounded $z$, the functions
\[
 h_N(z)=W_N(z/N)=\frac{(1-z/N)^N}{1+z/N}
\]
converge in $C^j$ to $e^{-z}$ for every fixed $j$. At $y=1$ use
\[
 \frac{h_N(cy^2)-h_N(c)}{1-y}
 =-c(1+y)\int_0^1 h_N'\bigl(cy^2+t c(1-y^2)\bigr)\dd t.
\]
This removes the apparent singularity and proves the stated uniformity.
\end{proof}

\subsection{The unique limiting maximum}

\begin{theorem}[Global kernel limit and maximizer localization]
\label{thm:kernel-global-limit}
There is exactly one $y_\infty\in(0,1)$ satisfying
\begin{equation}\label{eq:yinfty}
 1+y-y^2-y^3+4y\log y=0.
\end{equation}
Put
\begin{equation}\label{eq:cinfty}
 c_\infty=\frac{-2\log y_\infty}{1-y_\infty^2},\qquad
 M_\infty=(1+y_\infty)
      y_\infty^{2y_\infty^2/(1-y_\infty^2)}.
\end{equation}
Then
\begin{equation}\label{eq:Mnlimit}
 M_N\longrightarrow M_\infty>1.
\end{equation}
For \emph{every} choice of global maximizers $(r_N,y_N)$,
\begin{equation}\label{eq:kernel-localization}
 Nr_N\longrightarrow c_\infty,\qquad y_N\longrightarrow y_\infty.
\end{equation}
The limiting function $H$ has a unique global maximum at
$(c_\infty,y_\infty)$, with negative definite Hessian there.
\end{theorem}
\begin{proof}
For fixed $0<y<1$, direct differentiation shows that $H(c,y)$ has
exactly one maximum in $c$, attained at
$c(y)=-2\log y/(1-y^2)$; its value is $B(y)$.
The logarithmic derivative of $B$ has the sign of
\[
 F(y)=1+y-y^2-y^3+4y\log y,
 \qquad
 (\log B)'(y)=\frac{F(y)}{(1-y^2)^2}.
\]
We have $F(0+)=1$, $F(1)=F'(1)=0$, and
\[
 F''(y)=-\frac{2(3y-2)(y+1)}y.
\]
Thus $F'$ increases from $-\infty$ up to $y=2/3$, and then decreases
to zero. Its maximum is positive because
$F'(2/3)=7/3+4\log(2/3)>1/3$, using $\log(3/2)<1/2$.
It follows that $F$ first decreases and then increases, with exactly
one zero in $(0,1)$. At that zero $F'<0$, so $B''(y_\infty)<0$.
The limiting maximum exceeds one: at $y=1/4$,
$B(y)>1$ is equivalent to $(5/4)^{15}>16$.

The uniform envelope in \cref{lem:kernel-scale} gives
$\limsup M_N\le\max B=M_\infty$. Evaluating at
$(r,y)=(c_\infty/N,y_\infty)$ gives the matching lower bound.
If $(r_N,y_N)$ maximizes $K_N$, then
\[
 M_N\le B_N(y_N)=B(y_N)+o(1).
\]
Uniqueness of the maximum of $B$ and compactness force $y_N\to y_\infty$.
In particular $y_N$ stays away from zero and one. If $c_N=Nr_N$
were unbounded along a subsequence, the estimate
\[
 0\le K_N(r_N,y_N)
 \le\frac{e^{-Nr_Ny_N^2}}{1-y_N}
\]
would force its value to zero, a contradiction. Every bounded
subsequence has a further limit $c$, and the local scaling theorem
gives $H(c,y_\infty)=M_\infty$. The unique conditional maximum
then gives $c=c_\infty$; this proves \eqref{eq:kernel-localization}.

Finally $H_{cc}<0$ at each conditional maximum. Differentiating
$H(c(y),y)=B(y)$ shows that the Schur complement of $H_{cc}$ in
the full Hessian is
$H_{yy}-H_{cy}^2/H_{cc}=B''(y_\infty)<0$.
This proves negative definiteness.
\end{proof}

\subsection{Convergent reciprocal-index expansions}

\begin{theorem}[Eventual uniqueness and analytic perturbation]
\label{thm:kernel-expansion}
For all sufficiently large integers $N$, $K_N$ has exactly one
global maximizer. There are real analytic functions $c(s)$, $y(s)$,
and $M(s)$ near $s=0$ such that
\[
 Nr_N=c(1/N),\qquad y_N=y(1/N),\qquad M_N=M(1/N).
\]
Their Taylor series give convergent expansions in $1/N$ for large $N$.
In particular,
\begin{equation}\label{eq:Mn-first}
 M_N=M_\infty+\frac{\kappa_\infty}{N}+O(N^{-2}),\qquad
 \kappa_\infty=
 \frac{c_\infty^2y_\infty^2(1+y_\infty)
       e^{-c_\infty y_\infty^2}}2>0.
\end{equation}
Consequently $M_{N+1}<M_N$ for all sufficiently large $N$.
\end{theorem}
\begin{proof}
Define, with its removable value at $s=0$,
\[
 \mathcal W(s,z)=\exp\left(\frac{\log(1-sz)}s-\log(1+sz)\right),
 \qquad \mathcal W(0,z)=e^{-z},
\]
and
\[
 \mathcal K(s,c,y)=
 \frac{\mathcal W(s,cy^2)-\mathcal W(s,c)}{1-y}.
\]
It is analytic near $(0,c_\infty,y_\infty)$, and at $s=1/N$
it is $K_N(c/N,y)$. The analytic implicit function theorem applied
to $(\mathcal K_c,\mathcal K_y)$, whose Jacobian is the invertible
Hessian in \cref{thm:kernel-global-limit}, gives a unique analytic
stationary branch $(c(s),y(s))$. The localization theorem places
every global maximizer in an arbitrarily small neighborhood of this
point for large $N$. On a fixed smaller convex neighborhood, the Hessian
remains negative definite for small $s$. Strict concavity there and
the implicit-function uniqueness prove global uniqueness.

The first two terms are
\[
 \mathcal K(s,c,y)=H(c,y)+sH_1(c,y)+O(s^2),
\]
where, writing $A(z)=z+z^2/2$,
\begin{equation}\label{eq:H1}
 H_1(c,y)=\frac{e^{-c}A(c)-e^{-cy^2}A(cy^2)}{1-y}.
\end{equation}
Stationarity gives $e^{-c_\infty}=y_\infty^2e^{-c_\infty y_\infty^2}$.
Substitution into \eqref{eq:H1} gives exactly $\kappa_\infty$ in
\eqref{eq:Mn-first}. The first location correction is the explicit
matrix formula
\begin{equation}\label{eq:kernel-drift}
 \begin{pmatrix}c'(0)\\y'(0)\end{pmatrix}
 =-\bigl(D^2H(c_\infty,y_\infty)\bigr)^{-1}
        \nabla H_1(c_\infty,y_\infty).
\end{equation}
Higher corrections follow by ordinary power-series substitution.
Because the optimized value is itself analytic and
$M'(0)=\kappa_\infty>0$, it increases with $s$ near zero.
Taking $s=1/N$ proves eventual strict decrease. The two-term
asymptotic alone, without this derivative information, would not
justify the last conclusion.
\end{proof}

The diagnostic values, computed independently of the proof, are
\begin{align*}
 y_\infty&=0.1453038592593584238735558631\ldots,\\
 c_\infty&=3.9410648316472506681136817940\ldots,\\
 M_\infty&=1.0538619303627116260728728105\ldots,\\
 \kappa_\infty&=0.1727965939279004260875754084\ldots.
\end{align*}
The incoming conjecture concerning the limit and all maximizer sequences
is therefore settled. The stronger assertion of uniqueness and decrease
for \emph{every} $N\ge2$ remains a separate finite-index problem;
\cref{sec:exact-M2} resolves the uniqueness assertion at $N=2$.

